\section{Introduction}

Dementia affects about 57 million people worldwide, a number projected to rise to about 153 million by 2050
\citep{who2023dementia}; in India alone, cases are expected to grow from 5.1 million in 2021 to more than 20
million by 2050. Mild cognitive impairment (MCI) is the main prodromal stage of dementia \citep{petersen2004mci},
and it is the window in which intervention is most useful. Neuropsychological testing, magnetic resonance
imaging (MRI), positron emission tomography and cerebrospinal fluid biomarkers are costly, invasive or scarce
outside large hospitals. Scalp electroencephalography (EEG), in contrast, is cheap, portable, non-invasive and
widely available. Alzheimer's-type dementia produces well-documented EEG changes: slowing of the background
rhythm, reduced posterior alpha, increased theta and delta power, and reduced coherence
\citep{jeong2004eeg,dauwels2010diagnosis,babiloni2020measures}.

The CAUEEG dataset \citep{kim2023caueeg} is one of the largest public clinical EEG collections with a
Normal / MCI / Dementia benchmark. Results on it are, however, hard to interpret for four reasons. First,
clinical EEG slowing is read during eyes-closed rest, yet recordings also contain eyes-open periods, photic
(flashing-light) stimulation and hyperventilation, and random crops mix these states. Second, per-segment
scoring, or splits that share patients, inflate accuracy relative to per-recording, patient-disjoint
evaluation \citep{saeb2017voice}. Third, a 118-recording test set has a 95\% confidence interval of roughly
$\pm 9$ accuracy points \citep{varoquaux2018cross}. Fourth, deep models are rarely compared with a well-tuned
tabular model of established EEG markers under the same protocol, and rarely tested at another site. The
published models are also large (11--87 million parameters per model and 254 million for an ensemble), which
limits deployment on portable hardware.

\subsection{Objectives}
The objectives of this work are:
\begin{enumerate}[leftmargin=*,itemsep=0pt,topsep=2pt]
\item to analyse the CAUEEG dataset in depth and make its recording-state annotations usable;
\item to design a lightweight, interpretable deep learning model for three-class Normal / MCI / Dementia
      classification;
\item to compare it with tabular and deep learning baselines trained under exactly the same protocol;
\item to validate the model on independent datasets without retraining.
\end{enumerate}

\subsection{Contributions}
(1) A verified correction of CAUEEG event timing and a per-second state and artifact timeline.
(2) SANF-Net, a state-aware, age-normative fusion network of 143k parameters that uses alpha-reactivity and
photic-driving biomarkers, which can only be computed with the corrected clock. An ensemble of three SANF-Net
models reaches 72.9\% accuracy on the official test split.
(3) The identification of a photic-stimulation protocol confound between the Normal subgroups, separated from
the age confound by age matching.
(4) A patient-disjoint, per-recording evaluation with repeated five-fold cross-validation, bootstrap
confidence intervals, tabular and deep baselines, ablations, and external testing on two independent MCI
cohorts.
